\begin{ex}
Làm tròn số 43 đến hàng chục ta được số nào?
\shortans{40}
\loigiai{
Chữ số hàng đơn vị là 3 < 5 nên làm tròn thành 40.
}
\end{ex}

\begin{ex}
Làm tròn số 78 đến hàng chục ta được số nào?
\shortans{80}
\loigiai{
Chữ số hàng đơn vị là 8 $\ge$ 5 nên làm tròn thành 80.
}
\end{ex}

\begin{ex}
Làm tròn số 65 đến hàng chục ta được số nào?
\shortans{70}
\loigiai{
Chữ số hàng đơn vị là 5 nên làm tròn lên 70.
}
\end{ex}

\begin{ex}
Làm tròn số 234 đến hàng chục ta được số nào?
\shortans{230}
\loigiai{
Chữ số hàng đơn vị là 4 nên làm tròn thành 230.
}
\end{ex}

\begin{ex}
Làm tròn số 586 đến hàng chục ta được số nào?
\shortans{590}
\loigiai{
Chữ số hàng đơn vị là 6 nên làm tròn lên 590.
}
\end{ex}

\begin{ex}
Làm tròn số 315 đến hàng chục ta được số nào?
\shortans{320}
\loigiai{
Chữ số hàng đơn vị là 5 nên làm tròn lên 320.
}
\end{ex}

\begin{ex}
Làm tròn số 320 đến hàng trăm ta được số nào?
\shortans{300}
\loigiai{
Chữ số hàng chục là 2 < 5 nên làm tròn thành 300.
}
\end{ex}

\begin{ex}
Làm tròn số 670 đến hàng trăm ta được số nào?
\shortans{700}
\loigiai{
Chữ số hàng chục là 7 $\ge$ 5 nên làm tròn lên 700.
}
\end{ex}

\begin{ex}
Làm tròn số 450 đến hàng trăm ta được số nào?
\shortans{500}
\loigiai{
Chữ số hàng chục là 5 nên làm tròn lên 500.
}
\end{ex}

\begin{ex}
Làm tròn số 1248 đến hàng trăm ta được số nào?
\shortans{1200}
\loigiai{
Chữ số hàng chục là 4 < 5 nên làm tròn thành 1200.
}
\end{ex}

\begin{ex}
Làm tròn số 3682 đến hàng trăm ta được số nào?
\shortans{3700}
\loigiai{
Chữ số hàng chục là 8 $\ge$ 5 nên làm tròn lên 3700.
}
\end{ex}

\begin{ex}
Làm tròn số 5450 đến hàng trăm ta được số nào?
\shortans{5500}
\loigiai{
Chữ số hàng chục là 5 nên làm tròn lên 5500.
}
\end{ex}

\begin{ex}
Làm tròn số 7819 đến hàng chục ta được số nào?
\shortans{7820}
\loigiai{
Hàng đơn vị là 9 nên làm tròn lên 7820.
}
\end{ex}

\begin{ex}
Làm tròn số 9004 đến hàng chục ta được số nào?
\shortans{9000}
\loigiai{
Hàng đơn vị là 4 nên làm tròn xuống 9000.
}
\end{ex}

\begin{ex}
Làm tròn số 8975 đến hàng trăm ta được số nào?
\shortans{9000}
\loigiai{
Hàng chục là 7 nên làm tròn lên: $8900 + 100 = 9000$.
}
\end{ex}

\begin{ex}
Có bao nhiêu số tự nhiên có hai chữ số khi làm tròn đến hàng chục thì được số 50?
\shortans{10}
\loigiai{
Các số từ 45 đến 54 (gồm 10 số).
}
\end{ex}

\begin{ex}
Số nhỏ nhất khi làm tròn đến hàng chục được số 80 là số nào?
\shortans{75}
\loigiai{
Số nhỏ nhất là 75.
}
\end{ex}

\begin{ex}
Số lớn nhất khi làm tròn đến hàng chục được số 80 là số nào?
\shortans{84}
\loigiai{
Số lớn nhất là 84.
}
\end{ex}

\begin{ex}
Số nhỏ nhất khi làm tròn đến hàng trăm được số 600 là số nào?
\shortans{550}
\loigiai{
Số nhỏ nhất là 550.
}
\end{ex}

\begin{ex}
Số lớn nhất khi làm tròn đến hàng trăm được số 600 là số nào?
\shortans{649}
\loigiai{
Số lớn nhất là 649.
}
\end{ex}

\begin{ex}
Quãng đường từ nhà An đến trường dài 782 m. Làm tròn đến hàng trăm thì khoảng bao nhiêu mét?
\shortans{800}
\loigiai{
Làm tròn 782 đến hàng trăm được 800 m.
}
\end{ex}

\begin{ex}
Một cửa hàng bán được 148 quyển vở. Làm tròn đến hàng chục là bao nhiêu quyển vở?
\shortans{150}
\loigiai{
Làm tròn 148 đến hàng chục được 150.
}
\end{ex}

\begin{ex}
Cân nặng một con bê là 85 kg. Làm tròn đến hàng chục là bao nhiêu ki-lô-gam?
\shortans{90}
\loigiai{
Làm tròn lên thành 90 kg.
}
\end{ex}

\begin{ex}
Làm tròn số 349 đến hàng trăm ta được số nào?
\shortans{300}
\loigiai{
Hàng chục là 4 < 5 nên làm tròn thành 300.
}
\end{ex}

\begin{ex}
Làm tròn số 351 đến hàng trăm ta được số nào?
\shortans{400}
\loigiai{
Hàng chục là 5 nên làm tròn lên 400.
}
\end{ex}
